% Generated by analysis/s58_sm_tables.py from tab_extension.tex -- do not edit.
% The supplement needs its own label so that the main paper's "Supplementary
% Table" pointers resolve; everything else is the source table verbatim.
\begin{table}[t]
\centering
\caption{Forecasting the grid. On an ordinal basis the genotype model is a \emph{function} of a level rather than a lookup over levels, so it can be evaluated past the last level measured. Top: the test. An entire extreme level of one gene is removed from the fit---every architecture that uses it---and the model predicts the accuracy of those architectures; the error is compared against the same model inside the grid and against the spread of the held-out level itself. Bottom: the forecast, for each range whose optimum lies on its boundary \emph{on both datasets}, of the change in mean accuracy from one further level (90\% bootstrap interval over architectures). Ranges on the boundary of only one dataset are given in the text. Discarding the curvature term outside the range is the conservative reading of the same fit and is shown as a robustness check; across all four forecasts it changes no sign.}
\label{tab:sm_extension}
\footnotesize
\resizebox{\columnwidth}{!}{%
\begin{tabular}{lcc}
\toprule
 & CIFAR-10 & CIFAR-100 \\
\midrule
\multicolumn{3}{l}{\emph{Leave-one-level-out: a level the model never saw}} \\
\quad Error on the level's mean accuracy, mean / median (pts) & 2.89 / 1.12 & 1.94 / 1.69 \\
\quad Within 2 points & 7/10 & 6/10 \\
\quad $\rho$ inside the held-out level & 0.856 & 0.788 \\
\quad Per-arch.\ MAE, outside / inside the grid & 7.55 / 5.64 & 4.94 / 3.91 \\
\addlinespace[2pt]\midrule
\multicolumn{3}{l}{\emph{Forecast gain from one further level}} \\
\quad dropout\_prob: 0.05 $\to$ 0.025 & $+4.12$ [$+3.85$, $+4.44$] & $+2.71$ [$+2.50$, $+2.94$] \\
\quad \dots\ curvature discarded (tangent) & $+3.96$ & $+2.53$ \\
\bottomrule
\end{tabular}}
\end{table}
